\chapter{Estado del arte}\label{cap:estado-arte}

\section{Detección clásica de \emph{fake news}}

La línea inicial de trabajo en detección automática de noticias falsas se
apoya en representaciones léxicas y modelos lineales o ensembles. Las
\emph{features} habituales son TF-IDF sobre $n$-gramas de palabras y
caracteres \citep{conroy2015automatic}, complementadas con rasgos
estilométricos (longitud media de oración, riqueza léxica, índices de
legibilidad) y rasgos psicolingüísticos (LIWC). Sobre estas
representaciones se entrenan clasificadores como Naïve Bayes,
\emph{Support Vector Machines} (SVM) lineales con núcleo lineal o
\emph{gradient boosting}.

\citet{shu2017fake} ofrecen un panorama exhaustivo de la literatura
hasta esa fecha y subrayan la dificultad metodológica del campo: la
ausencia de un \emph{benchmark} consensuado, la dependencia del dominio
y la confusión recurrente entre veracidad fáctica, sesgo editorial y
sátira. Trabajos posteriores han mostrado que muchos modelos de alto
rendimiento sobre datasets cerrados no generalizan: la separación
\emph{real}/\emph{falso} se resuelve, en parte, identificando la
\emph{fuente} más que el contenido \citep{ahmed2017detection}.

\section{Inferencia de lenguaje natural y \textsc{FEVER}}

La inferencia de lenguaje natural (NLI), también llamada \emph{textual
entailment}, es la tarea de decidir, dados un par
$(\text{premisa}, \text{hipótesis})$, si la premisa implica
(\emph{entailment}), contradice (\emph{contradiction}) o es neutra
(\emph{neutral}) con respecto a la hipótesis. Datasets como SNLI
\citep{bowman2015snli} y MultiNLI \citep{williams2018multinli} han
permitido entrenar modelos profundos con cientos de miles de pares
anotados.

El dataset \textbf{FEVER} \citep{thorne2018fever} adapta esta
formulación a la verificación factual: las hipótesis son
\emph{claims} sintéticos generados por anotadores a partir de Wikipedia,
y las premisas son frases concretas extraídas como evidencia. Las tres
etiquetas son \textsc{supported}, \textsc{refuted} y
\textsc{not-enough-info}. La tarea completa de FEVER se descompone
canónicamente en tres subproblemas:

\begin{enumerate}[leftmargin=*,itemsep=0.2em]
  \item \emph{Document retrieval}: dado un \emph{claim}, recuperar las
        páginas de Wikipedia relevantes.
  \item \emph{Sentence selection}: dentro de esas páginas, identificar
        las frases que actúan como evidencia.
  \item \emph{Claim verification}: aplicar un clasificador NLI sobre los
        pares (claim, evidencia) y agregar para producir un veredicto.
\end{enumerate}

El presente trabajo se centra exclusivamente en el subproblema (3), el
único de naturaleza puramente estadística-predictiva; la recuperación se
sustituye por evidencias proporcionadas por una API de búsqueda web
descrita en la memoria de Informática.

\section{Modelos basados en \emph{transformers}}

A partir de \citet{vaswani2017attention} la arquitectura
\emph{transformer} desplaza progresivamente a redes recurrentes y
convolucionales en PLN. Modelos preentrenados como BERT
\citep{devlin2019bert}, RoBERTa \citep{liu2019roberta} y, más
recientemente, DeBERTa \citep{he2021deberta,he2023debertav3} han
establecido el estado del arte en NLI. DeBERTa-v3 introduce dos
mejoras relevantes para esta tarea: atención \emph{disentangled} (separa
contenido y posición) y un objetivo de preentrenamiento \emph{Replaced
Token Detection} más eficiente. Su versión \emph{base} (184\,M
parámetros) ofrece un compromiso favorable entre coste computacional y
precisión, y es la elegida para los experimentos de este trabajo.

\section{Calibración de probabilidades}

Aunque un clasificador alcance una precisión elevada, sus probabilidades
\emph{a posteriori} pueden estar sistemáticamente sesgadas, fenómeno
documentado de forma general en redes neuronales modernas
\citep{guo2017calibration}. La calibración \emph{post-hoc} mediante
\emph{Platt scaling} \citep{platt1999probabilistic} o regresión
isotónica \citep{zadrozny2002transforming} es una práctica
metodológicamente sólida cuando se persigue el uso de las probabilidades
como tales (por ejemplo, para fijar umbrales de decisión, integrarlas
en una regla de agregación, o reportarlas al usuario). El
\emph{Expected Calibration Error} (ECE) y los diagramas de fiabilidad
\citep{niculescu2005predicting} son las métricas estándar para
diagnóstico.

\section{Comparación estadística de clasificadores}

La comparación de dos clasificadores sobre un mismo conjunto de prueba
es un problema clásico de inferencia. \citet{dietterich1998approximate}
revisa los principales tests y recomienda, para muestras pareadas
moderadas, el test de McNemar sobre la matriz de discordancias. Para
estadísticos no diferenciables como el F1 macro, el \emph{bootstrap}
no paramétrico \citep{efron1993introduction} permite construir
intervalos de confianza válidos bajo supuestos mínimos. Estas técnicas
se aplican en el capítulo~\ref{cap:tests}.
